\documentclass[11pt]{article}
\usepackage[margin=0.9in]{geometry}
\usepackage[T1]{fontenc}
\setlength{\parskip}{0.6em}
\setlength{\parindent}{0pt}
\begin{document}
\section*{Integration and Support NIST Mapping}
The integration module coordinates the platform-level behavior that makes the individual service areas operate as one cloud system. In NIST terms, this module does not provide a separate essential characteristic by itself; instead, it connects rapid elasticity, resource pooling, and measured service into one usable operational surface. Health checks, startup synchronization, quota mode visibility, and unified project snapshots all contribute to platform manageability.

Phase 2 integration work focuses on end-to-end correctness. Scaling must respect quotas, monitoring data must remain readable through the project metrics API, and network cleanup must remain safe while projects grow or shrink. The integration view exposes these relationships through a single API payload so operators can confirm that the control path, runtime state, and reporting path are all aligned.

This module also supports service orchestration principles common in cloud systems by centralizing startup repair logic. Missing containers are marked appropriately, orphaned networks are cleaned up, and the operations console can show both platform-wide and project-specific state. That makes the overall data plane more reliable without adding a more complex distributed control mechanism.
\end{document}
